\clearpage
\thispagestyle{empty}
\refstepcounter{section}
\label{app:worked-example}

\begin{center}
\textbf{APPENDIX A}

\vspace{\baselineskip}

\textbf{Worked Example of the Generic GBBRPM\\
Core Architecture}
\end{center}

\vspace{3\baselineskip}

\subsection{Purpose and Scope of the Example}

This appendix provides a transparent, step-by-step calculation of the generic
Graph-Based Bounded Risk Propagation Model (GBBRPM) core architecture. The
example demonstrates source-node initialization, relation-level contribution,
bounded convergence, local disturbance at a receiving node, current-state
recursion, and final comparative ranking. All values are controlled,
dimensionless illustrative inputs. They are not drainage measurements,
software observations, electrical measurements, calibrated probabilities, or
evidence of empirical predictive validity.

The calculation uses the abstract directed acyclic graph shown in
Figure~\ref{fig:generic-node-example}. Its directed relations are
\(N_1\rightarrow N_2\), \(N_4\rightarrow N_2\), and
\(N_2\rightarrow N_3\). Nodes \(N_1\) and \(N_4\) are sources, the two
incoming paths converge at \(N_2\), and the accumulated state at \(N_2\) is
then propagated to \(N_3\). A valid topological evaluation order is
\(N_1,N_4,N_2,N_3\).

\subsection{Controlled Inputs}

Table~\ref{tab:appendix-node-inputs} lists the local disturbance assigned to
each node. The nonzero \(B_2\) value is included deliberately to show that a
receiving node can contain both locally introduced and inherited risk.

\begin{table}[H]
\caption{Controlled Node Inputs for the Worked Example}
\label{tab:appendix-node-inputs}
\centering
\begin{tabular*}{\textwidth}{@{\extracolsep{\fill}}lll}
\toprule
Node & Predecessor set \(P(j)\) & Local disturbance \(B_j\) \\
\midrule
\(N_1\) & \(\emptyset\) & 0.60 \\
\(N_4\) & \(\emptyset\) & 0.30 \\
\(N_2\) & \(\{N_1,N_4\}\) & 0.10 \\
\(N_3\) & \(\{N_2\}\) & 0.00 \\
\bottomrule
\end{tabular*}
\end{table}

Table~\ref{tab:appendix-edge-inputs} lists the relation-specific
susceptibilities and transmission factors. Neutral transmission
\(\tau_{ij}=1\) is used consistently with the frozen v1 baseline. Therefore,
the worked example does not claim to evaluate heterogeneous transmission.

\begin{table}[H]
\caption{Controlled Relation Inputs for the Worked Example}
\label{tab:appendix-edge-inputs}
\centering
\begin{tabularx}{\textwidth}{@{}c c c L@{}}
\toprule
Relation & Susceptibility \(S_{ij}\) & Transmission \(\tau_{ij}\) & Role \\
\midrule
\(N_1\rightarrow N_2\) & 0.50 & 1.00 & First input to convergence \\
\(N_4\rightarrow N_2\) & 0.40 & 1.00 & Second input to convergence \\
\(N_2\rightarrow N_3\) & 0.75 & 1.00 & Recursive downstream transfer \\
\bottomrule
\end{tabularx}
\end{table}

All supplied values lie in \([0,1]\), every relation has a valid endpoint,
and the graph contains no directed cycle. The inputs therefore satisfy the
declared bounded-input and DAG requirements before propagation begins.

\subsection{Core Update Rules}

For a directed relation \(i\rightarrow j\), the incoming contribution is
calculated as

\[
Q_{ij}=S_{ij}\tau_{ij}R_i.
\eqtag{A-1}
\]

where:

\noindent\(Q_{ij}\) = contribution entering node \(j\) from predecessor \(i\);\\
\noindent\(S_{ij}\) = bounded relation-specific susceptibility;\\
\noindent\(\tau_{ij}\) = bounded transmission or influence factor; and\\
\noindent\(R_i\) = current accumulated state of upstream node \(i\).

The receiving-node state is calculated as

\[
R_j=1-(1-B_j)\prod_{i\in P(j)}(1-Q_{ij}).
\eqtag{A-2}
\]

where:

\noindent\(R_j\) = final bounded state of receiving node \(j\);\\
\noindent\(B_j\) = local disturbance introduced directly at node \(j\);\\
\noindent\(P(j)\) = set of immediate predecessors of node \(j\); and\\
\noindent\(Q_{ij}\) = incoming contribution from predecessor \(i\).

The multiplicative-complement expression is a bounded saturating aggregator.
It is used algebraically and is not interpreted as a probabilistic union.

\subsection{Step 1: Initialize the Source Nodes}

Nodes \(N_1\) and \(N_4\) have no predecessors. For either source node, the
product over the empty predecessor set equals one. Equation A-2 therefore
reduces to \(R_j=B_j\):

\[
\begin{aligned}
R_1
  &=1-(1-B_1)\prod_{i\in\emptyset}(1-Q_{i1}) \\
  &=1-(1-0.60)(1) \\
  &=0.60,\\[0.5\baselineskip]
R_4
  &=1-(1-B_4)\prod_{i\in\emptyset}(1-Q_{i4}) \\
  &=1-(1-0.30)(1) \\
  &=0.30.
\end{aligned}
\eqtag{A-3}
\]

Thus, a source node can contain a local disturbance. It does not require a
special propagation rule: the ordinary aggregation rule produces its local
state automatically through the empty-product convention.

\subsection{Step 2: Calculate Contributions Entering
\texorpdfstring{\(N_2\)}{N2}}

The current source states are now available for the two relations entering
\(N_2\). Applying Equation A-1 gives

\[
\begin{aligned}
Q_{12}
  &=S_{12}\tau_{12}R_1 \\
  &=(0.50)(1.00)(0.60) \\
  &=0.30,\\[0.5\baselineskip]
Q_{42}
  &=S_{42}\tau_{42}R_4 \\
  &=(0.40)(1.00)(0.30) \\
  &=0.12.
\end{aligned}
\eqtag{A-4}
\]

The two contributions are calculated separately because they arrive through
different directed relations with different susceptibilities. Their raw sum
is not used as the receiving-node state; they are combined through the bounded
aggregation rule in the next step.

\subsection{Step 3: Aggregate Local and Inherited Risk at
\texorpdfstring{\(N_2\)}{N2}}

Node \(N_2\) has local disturbance \(B_2=0.10\) and two incoming
contributions. Substitution into Equation A-2 gives

\[
\begin{aligned}
R_2
  &=1-(1-B_2)(1-Q_{12})(1-Q_{42}) \\
  &=1-(1-0.10)(1-0.30)(1-0.12) \\
  &=1-(0.90)(0.70)(0.88) \\
  &=1-0.5544 \\
  &=0.4456.
\end{aligned}
\eqtag{A-5}
\]

The value \(R_2=0.4456\) represents the combined effect of the local
disturbance at \(N_2\) and the two inherited contributions. Although all
three components contribute to the result, the multiplicative-complement
form keeps the computed state within \([0,1]\).

For comparison, removing the local disturbance while retaining both inherited
contributions would give
\(1-(1-0.30)(1-0.12)=0.3840\). The increase from 0.3840 to 0.4456 is therefore
caused by including the receiving node's own local disturbance through the
same bounded rule rather than adding it after propagation.

\subsection{Step 4: Propagate the Current State to
\texorpdfstring{\(N_3\)}{N3}}

The contribution from \(N_2\) to \(N_3\) uses the newly accumulated state
\(R_2=0.4456\), not only the original local disturbance \(B_2=0.10\):

\[
\begin{aligned}
Q_{23}
  &=S_{23}\tau_{23}R_2 \\
  &=(0.75)(1.00)(0.4456) \\
  &=0.3342.
\end{aligned}
\eqtag{A-6}
\]

Because \(N_3\) has no local disturbance and only \(N_2\) as its
predecessor, its state is

\[
\begin{aligned}
R_3
  &=1-(1-B_3)(1-Q_{23}) \\
  &=1-(1-0.00)(1-0.3342) \\
  &=1-(1.00)(0.6658) \\
  &=0.3342.
\end{aligned}
\eqtag{A-7}
\]

This step demonstrates current-state recursion. If the calculation incorrectly
used only \(B_2\), the relation contribution would be
\((0.75)(1.00)(0.10)=0.0750\), which would discard the risk inherited by
\(N_2\). The implemented architecture instead propagates \(R_2\), preserving
the accumulated upstream effect in the downstream evaluation.

\subsection{Step 5: Produce the Comparative Ranking}

After all nodes have been evaluated once in topological order, their final
states are sorted in descending order. Table~\ref{tab:appendix-final-ranking}
reports the resulting prioritization.

\begin{table}[H]
\caption{Final Node States and Comparative Ranking}
\label{tab:appendix-final-ranking}
\centering
\begin{tabular*}{0.72\textwidth}{@{\extracolsep{\fill}}cll}
\toprule
Rank & Node & Final state \(R_i\) \\
\midrule
1 & \(N_1\) & 0.6000 \\
2 & \(N_2\) & 0.4456 \\
3 & \(N_3\) & 0.3342 \\
4 & \(N_4\) & 0.3000 \\
\bottomrule
\end{tabular*}
\end{table}

The ranking is a comparison among modeled node states under this controlled
configuration. It does not mean that \(N_1\) has a 60\% probability of failure
or that the numerical spacing between ranks has a domain-independent physical
meaning. If two nodes obtain exactly equal final states, the implementation
retains the tie and assigns them the same risk standing rather than applying an
unsupported secondary tie-breaker.

\subsection{Calculation Trace and Interpretation}

Table~\ref{tab:appendix-trace} consolidates the complete evaluation trace. It
also shows that every downstream calculation uses only predecessor states that
have already been computed in the declared topological order.

\begin{table}[H]
\caption{Complete Topological Evaluation Trace}
\label{tab:appendix-trace}
\centering
\small
\begin{tabularx}{\textwidth}{@{}c c L L c@{}}
\toprule
Order & Node & Available predecessor states & Incoming contributions & Result \\
\midrule
1 & \(N_1\) & None & Empty product & \(R_1=0.6000\) \\
2 & \(N_4\) & None & Empty product & \(R_4=0.3000\) \\
3 & \(N_2\) & \(R_1=0.6000\), \(R_4=0.3000\) & \(Q_{12}=0.3000\), \(Q_{42}=0.1200\) & \(R_2=0.4456\) \\
4 & \(N_3\) & \(R_2=0.4456\) & \(Q_{23}=0.3342\) & \(R_3=0.3342\) \\
\bottomrule
\end{tabularx}
\end{table}

The worked example establishes seven features of the generic architecture.
First, the empty-product convention gives source nodes their local states
without a separate initialization equation. Second, relation susceptibility
modifies each upstream state before it enters the receiving node. Third,
neutral transmission preserves the transmission interface without introducing
unsupported relation differences. Fourth, multiple incoming contributions and
a receiving-node disturbance are combined through one bounded rule. Fifth,
the evaluator propagates the current accumulated state rather than only the
original local input. Sixth, one topological pass is sufficient because the
evaluated graph is acyclic. Seventh, prioritization is derived only after all
node states have been calculated.

These calculations explain the mechanics of the frozen core architecture.
They do not establish that the illustrative parameter values are appropriate
for a particular domain. Drainage, software-dependency, and electrical-network
instantiations require separate evidence for the meanings and construction of
their local disturbances, susceptibilities, transmission factors, and outcome
comparators.
